1. Erdman — Introduction to Lie Algebras
2. Humphreys — Introduction to Lie Algebras and Representation Theory
3. Fulton — Representation theory
4. Goodman — Symmetry, Representations, and Invariants

1. Start with Erdman. Just until the proof of irreducible representation of SL2
2. Then follow with Humphreys which is mainly language.
3. Then, Fulton which is a classic,
4. and follow by Goodman which is modern.



Short wiki quotes on 9 authors of 8 mentioned introductiory texts on
representation theory:

Alexander Kirillov Jr. (?) is a mathematician working in the area of
representation theory and Lie groups. (cf. Representation Theory (Lie algebras,
quantum groups, affine Lie algebras, Conformal Field Theory) on his homepage)

Pavel Etingof (1969) does research on the intersection of mathematical physics
(exactly integrable systems) and representation theory, e.g. quantum groups.

Daniel Bump (1952) is a mathematician with contributions to number theory,
representation theory, combinatorics, and random matrix theory, as well as
mathematical exposition.

Joe Harris (1951) is an algebraic geometer.

Barry Simon (1946) is a mathematical physicist known for contributions in
spectral theory, functional analysis, and nonrelativistic quantum mechanics
(particularly Schrödinger operators), including the connections to atomic and
molecular physics.

James E. Humphreys (1939-2020) was a mathematician who worked in algebraic
groups, Lie groups, and Lie algebras and applications of these mathematical
structures.

William Fulton (1939) is an algebraic geometer.

Willard Miller Jr. (1937) is a mathematician with research in superintegrable
systems, Lie groups and algebras, special functions, q-series, mathematical
physics.

Alexandre Kirillov (1936) is a mathematician known for his works in the fields
of representation theory, topological groups and Lie groups. In particular he
introduced the orbit method into representation theory.


##
A summary of bibliography, lectures notes and software.

Basic Bibliography:
1) Simon: Representations of Finite and Compact Groups, 1995, 280 pp.
2) Kirillov Jr: An Introduction to Lie Groups and Lie Algebras, 2008, 177 / 230 pp. [ no finite groups ]
3) Humphreys: Introduction to Lie Algebras and Representation Theory, 1973, 186 pp. [ no finite groups ]

Additional Bibliography:
1') Fulton-Harris: Representation Theory, 1991, 560 pp.
2') Bump: Lie Groups, 2011, 532 pp.
3') Miller: Symmetry Groups and their Applications, 1972, 436 pp.

Bonus:
A) Kirillov Sr: Introduction to the Theory of Representations and Noncommutative Harmonic Analysis, 1988, 158 pp.
B) Khovanov: https://www.math.columbia.edu/~khovanov/finite (many links below can be found there)
C) Etingof: Introduction to Representation Theory, 2016, 231 pp.
D) Lam: Representations of Finite Groups: A Hundred Years, 1998. Part I (12 pp.), Part II (10 pp.)
E) Teleman: Representation Theory, 2005, 60 pp. (lecture notes)
F) Webb: A Course in Finite Group Representation Theory, 2016, 294 pp.
G) Baker: Representations of Finite Groups, 2019, 77 pp.
H) Jackson: Notes on the representation theory of finite groups, 2004, 55pp.

Software:
I. SAGE (some software below as GAP is bundles in SAGE also). In particular, see documentation Lie Methods and Related Combinatorics in Sage, by Bump, Salisbury, Schilling.
II. LiE
III. GAP (Groups, Algorithms, Programming - a System for Computational Discrete Algebra). See also Representations and Characters of Groups review of GAP's mathematical capabilities. Also teaching materials and software packages by Webb.
